\chapter{Conclusion and Future Work}
\label{ch:conclusion}

\section{Conclusions}

This work successfully designed and deployed an autonomous, end-to-end DevSecOps pipeline
capable of intelligently triaging static analysis alerts and self-healing vulnerabilities
without human intervention. The key findings are:

\begin{itemize}
  \item \textbf{Comprehensive SAST Integration:} Automated deep static application
        security testing using CodeQL as the foundational alert generator, producing
        SARIF reports across 2,766 OWASP Benchmark Java test cases with high recall
        but unacceptably low precision ($\approx$60\%), motivating the ML triage layer.

  \item \textbf{High-Fidelity Machine Learning Triage:} Upgraded from basic SARIF
        metadata features to an optimized Gradient Boosting Machine (GBM) utilizing
        13-dimensional weighted vulnerability vectors. GBM achieved the best overall
        performance with an F1-Score of 0.8887 and ROC-AUC of 0.9106, outperforming
        Random Forest and SVM classifiers.

  \item \textbf{Elimination of Alert Fatigue:} By mathematically optimizing the
        decision threshold (RF: 0.414, GBM: 0.364), the ML model suppressed static
        analysis noise and drastically improved triage precision to $\approx$85\%
        while retaining 93\% recall, ensuring critical vulnerabilities are not missed.

  \item \textbf{Autonomous LLM Remediation:} Engineered a sophisticated self-healing
        agent (Llama-3.3-70b via Groq API) that dynamically scales between a
        localized Snippet-Mode ($\pm$5-line context) and Full-File Escalation to
        correctly patch complex vulnerabilities including SQL Injection, XSS, and
        Path Traversal.

  \item \textbf{Zero-Trust Validation Loop:} Guaranteed strict codebase stability by
        forcing all LLM-generated patches to pass both a Java compiler check
        (\texttt{mvn clean compile}) and a CodeQL security re-scan before being
        automatically committed to the repository.
\end{itemize}

The combined result supports the thesis that a layered pipeline---combining static
analysis for coverage, ML classification for precision, and agentic AI for
actionability---is significantly more effective than any single component alone.

\section{Future Work}

Future enhancements to the pipeline include:
\begin{enumerate}
  \item \textbf{DAST Integration:} Incorporate OWASP ZAP to test running applications
        in real-time, catching runtime vulnerabilities and authentication flaws that
        are missed by SAST alone.

  \item \textbf{Multi-Language Expansion:} Extend support beyond Java to Python,
        JavaScript, and Go, creating a universal security layer for polyglot
        microservices architectures.

  \item \textbf{Reinforcement Learning with Compiler Feedback (RLCF):} Implement a
        feedback loop to fine-tune the LLM (Llama-3) based on compile success and
        CodeQL re-scan outcomes, progressively improving zero-shot patching accuracy
        over time.

  \item \textbf{Autonomous Test Generation:} Upgrade the agentic layer to
        automatically generate unit test cases alongside security patches, preventing
        future vulnerability regressions in patched components.
\end{enumerate}
